\section{Related Work}
% provenance: jrn_01KXT4CF02R1VRTMK9M3XQ6SEC (Passos differentiation); dec_01KXT28WTZ92PJ7S8D3MME78D4 (novelty framing).
% All citations verified against Crossref/OpenAlex/DBLP/IETF/publisher; no placeholders remain.

Table~\ref{tab:relatedcmp} compares CONTAINDER against the concrete systems and standards that
answer parts of the same problem, rather than against categories of work.

\emph{Workload identity is the closest existing system.} SPIFFE and its SPIRE
runtime~\cite{spiffe-spec,spiffe-cncf} implement almost exactly the chain we describe: attest node
and workload, issue short-lived scope-bound X.509 SVIDs keyed to the attested identity, refuse
renewal when attestation fails, rely on expiry rather than revocation. Recent work extends it to
IoT fleets~\cite{bakoyiannis2025spiffe} and nested tokens~\cite{cochak2024lsvid}. The question is
therefore whether CONTAINDER is workload identity with a DER label, and two differences are
load-bearing. First, SPIFFE has no notion of a \emph{physical effect that outlives the
credential}: withdrawing a workload identity stops a process opening new connections, but it does
not retract a volt-var curve already written into an inverter. Our $T_{\mathrm{sess}}$ and
$T_{\mathrm{cmd}}$ enforcement exists because actuation persists in a way computation does not,
and Section~\ref{sec:results} measures how far each layer reaches. Second, scope binds to
\texttt{FunctionSetAssignments}, whose conformant settings are dictated by
IEEE~1547-2018~\cite{ieee1547-2018} and CSIP~\cite{sunspec2018csip} rather than chosen freely.

\emph{The Web PKI has already adopted non-renewal.} The CA/Browser Forum defined a Short-lived
Subscriber Certificate in 2023 and \emph{exempted it from CRL and OCSP status
checking}~\cite{cabf-sc063v4} --- ``non-renewal replaces revocation'', already ratified, with
Baseline Requirements v2.2.8 setting the threshold at seven days~\cite{cabf-br} --- and ballot
SC-081v3 phases maximum TLS validity from 398 days to 47 by 2029~\cite{cabf-sc081v3}. This is at
once the strongest external validation of the premise and the sharpest limit on our novelty. We
claim none for short-lived credentials or for expiry-based containment. The same limit applies
inside our own problem domain: the Sandia analysis of the DER standards recommends limited-life
certificates and states that manufacturers following the IEEE~2030.5 guidance are ``not required
to limit the life of a device certificate''~\cite{obert2019recommendations}, which makes a
short-lived operational credential \emph{permitted} rather than novel. What we claim is the
mapping onto DER authorization: which scope object to bind, which control point to withhold, and
that session and command-effect enforcement are separately necessary because grid actuation
persists.

\emph{Competing key management for power systems.} IEC~62351 is the sharpest comparison: Part~8
gives role-based access control~\cite{iec62351-8} and Part~9 specifies key management
\emph{including} enrollment and revocation for power-system equipment~\cite{iec62351-9}. Part~9 is
a directly competing answer, and the honest question is why 2030.5/CSIP does not adopt it. Not
technical superiority on our side: 62351-9's revocation machinery assumes the connectivity and the
responsible reporting party that the Sandia analysis found
absent~\cite{obert2019recommendations}, and CSIP deployments certify against the 2030.5 profile.
Where an operator has that connectivity, 62351-9 is a reasonable alternative, and we say so rather
than claim the field. IEEE~1547.3-2023 is the cybersecurity guide for interconnected
DER~\cite{ieee1547.3-2023} and the natural home for these recommendations;
NISTIR~7628~\cite{nistir7628r1} and IEC~62443~\cite{iec62443-3-3,iec62443-4-2} give the enclosing
frameworks; NREL's DER certification line~\cite{nrel80581,saleem2022dercert} and gap
analysis~\cite{magill2026dergap} document the certification landscape.

\emph{IEEE~2030.5 and CSIP security.} The 2018 and 2023 editions mandate mutual TLS with X.509
device certificates~\cite{ieee20305-2018,ieee20305-2023}, constrained for interconnection by
CSIP~\cite{sunspec2018csip}. The closest prior work is the 2030.5 security tutorial of Passos et
al.~\cite{passos2025tutorial}, a survey that proposes no redesign, quantifies no feeder
consequence, and formalizes no impact metric; Qi et al.\ survey the broader DER and smart-inverter
surface~\cite{qi2016cybersecurity}.

\emph{DER and smart-inverter attacks.} High-wattage IoT botnets can move the
grid~\cite{soltan2018blackiot,shekari2022madiot}; inverter attacks destabilize grid and
market~\cite{hui2025destabilizing}; distribution voltage control is attackable under high
PV~\cite{isozaki2016detection}; and the SUN:DOWN disclosure catalogs fleet-scale inverter
vulnerabilities~\cite{forescout2025sundown}. We tie the attack surface to the credential lifecycle
rather than assuming device compromise. Distribution-grid consequence is quantified by the
false-data-injection literature~\cite{liu2009false,deng2017false,liang2017review}, grounded by the
2015 Ukraine incident~\cite{liang2017ukraine} and quantitative risk
formulations~\cite{teixeira2015secure}; our feeder method is supporting, not novel.

\emph{Building blocks.} ACME~\cite{rfc8555} and Let's Encrypt~\cite{aas2019letsencrypt} automate
issuance, SCEP~\cite{rfc8894} is the incumbent device enrollment, and HTTPS-ecosystem measurement
quantifies the lifecycle problem~\cite{durumeric2013analysis}. Secure device
identity~\cite{ieee8021ar2018}, firmware TPM~\cite{raj2016ftpm} and hardware
DICE~\cite{jager2020dice} root the bootstrap layer; the RATS architecture~\cite{rfc9334}, the
Entity Attestation Token~\cite{rfc9711} and BRSKI~\cite{rfc8995} carry attestation, on established
principles~\cite{coker2011principles}. Web-PKI revocation delivery is
gap-ridden~\cite{liu2015endtoend}, scalable revocation remains
open~\cite{larisch2017crlite,smith2020letsrevoke}, and OCSP Must-Staple adoption is
weak~\cite{chung2018ocsp}. Attack-graph metrics~\cite{ou2005mulval,ammann2002scalable,
jajodia2005topological,wang2008attackgraph,frigault2008measuring,poolsappasit2012dynamic,
wang2014kzero} inform the notation of Section~\ref{sec:formalism}, which we do not claim as a
contribution. On the power side, IEEE~1547~\cite{ieee1547-2018} defines the DER functions, DERMS
perform feeder voltage management~\cite{rahman2021derms}, role-based access
control~\cite{sandhu1996rbac} abstracts the scoping we narrow, and HELICS~\cite{hardy2024helics,
palmintier2017helics} and GridLAB-D~\cite{chassin2008gridlabd} are the co-simulation tooling a
cross-solver check would use.

\begin{table*}[t]
\centering
\caption{Positioning against the concrete systems and standards that answer part of the same
problem. \emph{Short-lived} = credential lifetime is the containment mechanism; \emph{Attested} =
issuance is gated on fresh evidence about the endpoint; \emph{Numeric scope} = the authorization
can bound the \emph{value} a control may take, not only which function may be invoked;
\emph{Session expiry} = an established session stops accepting commands when the credential
authorizing it expires; \emph{Command cleanup} = an already-issued control is retracted;
\emph{Physical eval.} = consequence is evaluated on a power-flow model.
y = yes, n = no, p = partial, --~= not applicable.}
\label{tab:relatedcmp}
\footnotesize
\begin{tabular}{@{}p{0.27\textwidth}ccccccc@{}}
\toprule
System or standard & DER/2030.5 & Short-lived & Attested & Numeric scope & Session expiry &
Command cleanup & Physical eval. \\
\midrule
IEEE 2030.5 / CSIP baseline~\cite{ieee20305-2018,sunspec2018csip} & y & n$^{\dagger}$ & n & n & n & n & n \\
SPIFFE / SPIRE~\cite{spiffe-spec,spiffe-cncf} & n & y & y & n & p$^{\ddagger}$ & -- & n \\
BRSKI~\cite{rfc8995} & n & p & y & n & n & -- & n \\
CA/B Forum short-lived certs~\cite{cabf-sc063v4,cabf-br} & n & y & n & n & n & -- & n \\
IEC 62351-8 (RBAC)~\cite{iec62351-8} & p & n & n & p & n & n & n \\
IEC 62351-9 (key mgmt)~\cite{iec62351-9} & p & p & n & n & n & n & n \\
\textbf{CONTAINDER (this work)} & y & y & p$^{\S}$ & y & y & y & y \\
\bottomrule
\end{tabular}
\\[2pt]
{\footnotesize
$^{\dagger}$Device certificates have indefinite lifetime and CRL/OCSP are prohibited for
IEEE~2030.5 certificates; local allow/deny listing is permitted but not
required~\cite{obert2019recommendations}.
$^{\ddagger}$SVID rotation bounds new connections; SPIFFE defines no retraction of an effect
already actuated, because a workload has none.
$^{\S}$Partial: the attestation path is implemented and tested but software-emulated. No
hardware-backed key or attestation-quote measurement is reported (Section~\ref{sec:limits}).}
\end{table*}
